\begin{thebibliography}{99}
\addcontentsline{toc}{chapter}{Referencias}
\setlength{\itemsep}{0.3em}

\bibitem[Advani et~al.(2019)]{advani2019}
Advani, A., Kitagawa, T., \& Słoczyński, T. (2019). Mostly harmless simulations? Using Monte Carlo studies for estimator selection. \textit{Journal of Applied Econometrics, 34}(6), 893--910. https://doi.org/10.1002/jae.2724

\bibitem[Athey et~al.(2019)]{athey2019}
Athey, S., Tibshirani, J., \& Wager, S. (2019). Generalized random forests. \textit{The Annals of Statistics, 47}(2), 1148--1178. https://doi.org/10.1214/18-AOS1709

\bibitem[Bach et~al.(2022)]{bach2022}
Bach, P., Chernozhukov, V., Kurz, M. S., \& Spindler, M. (2022). DoubleML: An object-oriented implementation of double machine learning in Python. \textit{Journal of Machine Learning Research, 23}(53), 1--6. https://www.jmlr.org/papers/v23/21-0862.html

\bibitem[Barrientos et~al.(2026)]{barrientos2026}
Barrientos, R., Delgado, R., Escalante, L., Febres, G., Hisbes, B., Chavez-Bedoya, L., Rosales, F., \& Vargas, W. (2026). Using machine learning to predict apartment prices in Lima. En \textit{Entrepreneurship and human-centric business strategies for social and economic resilience} (pp. 2725--2737). Springer. https://doi.org/10.1007/978-981-95-6415-6\_168

\bibitem[Boeing y Waddell(2017)]{boeing2017}
Boeing, G., \& Waddell, P. (2017). New insights into rental housing markets across the United States: Web scraping and analyzing Craigslist rental listings. \textit{Journal of Planning Education and Research, 37}(4), 457--476. https://doi.org/10.1177/0739456X16664789

\bibitem[Black(1999)]{black1999}
Black, S. E. (1999). Do better schools matter? Parental valuation of elementary education. \textit{The Quarterly Journal of Economics, 114}(2), 577--599. https://doi.org/10.1162/003355399556070

\bibitem[Cacñahuaray y Jiménez(2024)]{cacnahuaray2024}
Cacñahuaray Huaraca, J. C., \& Jiménez Novoa, S. M. (2024). \textit{Modelo predictivo de precios de alquiler de departamentos en la provincia de Lima usando técnicas de machine learning} [Tesis de segunda especialidad, Universidad Peruana Unión]. Repositorio institucional UPeU.

\bibitem[Cameron y Miller(2015)]{cameron2015}
Cameron, A. C., \& Miller, D. L. (2015). A practitioner's guide to cluster-robust inference. \textit{Journal of Human Resources, 50}(2), 317--372. https://doi.org/10.3368/jhr.50.2.317

\bibitem[Chay y Greenstone(2005)]{chay2005}
Chay, K. Y., \& Greenstone, M. (2005). Does air quality matter? Evidence from the housing market. \textit{Journal of Political Economy, 113}(2), 376--424. https://doi.org/10.1086/427462

\bibitem[Chen et~al.(2026)]{chen2026}
Chen, Z., Moeltner, K., \& Mair, M. (2026). \textit{Recovering direct price effects of environmental amenities in housing markets: Regression and causal machine learning model assessment with empirical Monte Carlo simulation} (arXiv:2606.02795). arXiv. https://arxiv.org/abs/2606.02795

\bibitem[Chernozhukov et~al.(2018)]{chernozhukov2018}
Chernozhukov, V., Chetverikov, D., Demirer, M., Duflo, E., Hansen, C., Newey, W., \& Robins, J. (2018). Double/debiased machine learning for treatment and structural parameters. \textit{The Econometrics Journal, 21}(1), C1--C68. https://doi.org/10.1111/ectj.12097

\bibitem[Chernozhukov et~al.(2022)]{chernozhukov2022}
Chernozhukov, V., Cinelli, C., Newey, W., Sharma, A., \& Syrgkanis, V. (2022). \textit{Long story short: Omitted variable bias in causal machine learning} (NBER Working Paper No. 30302). National Bureau of Economic Research. https://www.nber.org/papers/w30302

\bibitem[Cinelli et~al.(2024)]{cinelli2024}
Cinelli, C., Forney, A., \& Pearl, J. (2024). A crash course in good and bad controls. \textit{Sociological Methods \& Research, 53}(3), 1071--1104. https://doi.org/10.1177/00491241221099552

\bibitem[Cinelli y Hazlett(2020)]{cinelli2020}
Cinelli, C., \& Hazlett, C. (2020). Making sense of sensitivity: Extending omitted variable bias. \textit{Journal of the Royal Statistical Society: Series B, 82}(1), 39--67. https://doi.org/10.1111/rssb.12348

\bibitem[Contreras et~al.(2024)]{contreras2024}
Contreras, C., Aguilar, D., \& Lázaro, J. (2024). Impacto del transporte masivo en la informalidad de vivienda: el caso de la Línea 1 del Metro de Lima. \textit{Argumentos, 5}(1). Instituto de Estudios Peruanos. https://doi.org/10.46476/ra.v5i1.183

\bibitem[Egami et~al.(2023)]{egami2023}
Egami, N., Hinck, M., Stewart, B. M., \& Wei, H. (2023). Using imperfect surrogates for downstream inference: Design-based supervised learning for social science applications of large language models. En \textit{Advances in Neural Information Processing Systems 36} (NeurIPS 2023).

\bibitem[Faxvaag et~al.(2026)]{faxvaag2026}
Faxvaag, H., Aarsnes, U., \& Oust, A. (2026). Incorporating large language models in automated real estate valuation. \textit{Journal of Real Estate Research, 48}(3), 454--479. https://doi.org/10.1080/08965803.2025.2587313

\bibitem[Gabbe y Pierce(2017)]{gabbe2017}
Gabbe, C. J., \& Pierce, G. (2017). Hidden costs and deadweight losses: Bundled parking and residential rents in the metropolitan United States. \textit{Housing Policy Debate, 27}(2), 217--229. https://doi.org/10.1080/10511482.2016.1205647

\bibitem[Gyger et~al.(2026)]{gyger2026}
Gyger, T., Hauri, S., Bühlmann, S., Lehner, M., Schlesinger, J., \& Sigrist, F. (2026). Explainable spatial machine learning for hedonic real estate modeling. \textit{Real Estate Economics, 54}(3), 761--806. https://doi.org/10.1111/1540-6229.70030

\bibitem[Hernán y Robins(2020)]{hernan2020}
Hernán, M. A., \& Robins, J. M. (2020). \textit{Causal inference: What if}. Chapman \& Hall/CRC.

\bibitem[Hünermund et~al.(2023)]{hunermund2023}
Hünermund, P., Louw, B., \& Caspi, I. (2023). Double machine learning and automated confounder selection: A cautionary tale. \textit{Journal of Causal Inference, 11}(1), 20220078. https://doi.org/10.1515/jci-2022-0078

\bibitem[Imbens y Rubin(2015)]{imbens2015}
Imbens, G. W., \& Rubin, D. B. (2015). \textit{Causal inference for statistics, social, and biomedical sciences: An introduction}. Cambridge University Press.

\bibitem[Keele y Titiunik(2015)]{keele2015}
Keele, L. J., \& Titiunik, R. (2015). Geographic boundaries as regression discontinuities. \textit{Political Analysis, 23}(1), 127--155. https://doi.org/10.1093/pan/mpu014

\bibitem[Knaus et~al.(2021)]{knaus2021}
Knaus, M. C., Lechner, M., \& Strittmatter, A. (2021). Machine learning estimation of heterogeneous causal effects: Empirical Monte Carlo evidence. \textit{The Econometrics Journal, 24}(1), 134--161. https://doi.org/10.1093/ectj/utaa014

\bibitem[Lima Cómo Vamos(2016)]{limacomovamos2016}
Lima Cómo Vamos. (2016). \textit{Cómo vamos en movilidad: sexto informe de resultados sobre calidad de vida en Lima y Callao 2015}. Lima Cómo Vamos.

\bibitem[Manville(2013)]{manville2013}
Manville, M. (2013). Parking requirements and housing development: Regulation and reform in Los Angeles. \textit{Journal of the American Planning Association, 79}(1), 49--66. https://doi.org/10.1080/01944363.2013.785346

\bibitem[Mayo(2017)]{mayo2017}
Mayo Urtecho, H. M. (2017). \textit{Elaboración de un índice de precios de viviendas nuevas para Lima Metropolitana mediante precios hedónicos} [Tesis de grado, Universidad Nacional de Ingeniería]. http://hdl.handle.net/20.500.14076/13730

\bibitem[Municipalidad de Miraflores(2011)]{miraflores2011}
Municipalidad de Miraflores. (2011). \textit{Ordenanza N.° 342/MM: aprueban parámetros urbanísticos y edificatorios y las condiciones generales de edificación en el distrito} (art. 10: número mínimo de estacionamientos por unidad de vivienda; modificada por las Ordenanzas 385-MM y 466-MM).

\bibitem[Municipalidad de San Isidro(2020)]{sanisidro2020}
Municipalidad de San Isidro. (2020). \textit{Ordenanza N.° 523-MSI} (parámetros de estacionamiento para edificaciones multifamiliares).

\bibitem[Malpezzi(2003)]{malpezzi2003}
Malpezzi, S. (2003). Hedonic pricing models: A selective and applied review. En T. O'Sullivan \& K. Gibb (Eds.), \textit{Housing economics and public policy} (pp. 67--89). Blackwell.

\bibitem[Moran(1950)]{moran1950}
Moran, P. A. P. (1950). Notes on continuous stochastic phenomena. \textit{Biometrika, 37}(1/2), 17--23. https://doi.org/10.2307/2332142

\bibitem[Mundaca y Sánchez(2018)]{mundaca2018}
Mundaca, F., \& Sánchez, E. (2018). Índice de precios de inmuebles: un enfoque hedónico. \textit{Revista Estudios Económicos, 36}, 55--74. Banco Central de Reserva del Perú.

\bibitem[Nowak y Smith(2017)]{nowak2017}
Nowak, A., \& Smith, P. (2017). Textual analysis in real estate. \textit{Journal of Applied Econometrics, 32}(4), 896--918. https://doi.org/10.1002/jae.2550

\bibitem[Pearl(2009)]{pearl2009}
Pearl, J. (2009). \textit{Causality: Models, reasoning, and inference} (2.ª ed.). Cambridge University Press.

\bibitem[Pedregosa et~al.(2011)]{pedregosa2011}
Pedregosa, F., Varoquaux, G., Gramfort, A., Michel, V., Thirion, B., Grisel, O., et al. (2011). Scikit-learn: Machine learning in Python. \textit{Journal of Machine Learning Research, 12}, 2825--2830.

\bibitem[Rivera(2018)]{rivera2018}
Rivera Romero, L. A. (2018). \textit{Caso: determinantes del precio por metro cuadrado para departamentos en Lima Metropolitana: período 2005--2017} [Trabajo académico, Universidad de Lima].

\bibitem[Robinson(1988)]{robinson1988}
Robinson, P. M. (1988). Root-N-consistent semiparametric regression. \textit{Econometrica, 56}(4), 931--954. https://doi.org/10.2307/1912705

\bibitem[Rosen(1974)]{rosen1974}
Rosen, S. (1974). Hedonic prices and implicit markets: Product differentiation in pure competition. \textit{Journal of Political Economy, 82}(1), 34--55. https://doi.org/10.1086/260169

\bibitem[Shen y Ross(2021)]{shen2021}
Shen, L., \& Ross, S. (2021). Information value of property description: A machine learning approach. \textit{Journal of Urban Economics, 121}, 103299. https://doi.org/10.1016/j.jue.2020.103299

\bibitem[Sirmans et~al.(2005)]{sirmans2005}
Sirmans, G. S., Macpherson, D. A., \& Zietz, E. N. (2005). The composition of hedonic pricing models. \textit{Journal of Real Estate Literature, 13}(1), 1--44. https://doi.org/10.1080/10835547.2005.12090154

\bibitem[Tapia et~al.(2025)]{tapia2025}
Tapia, J., Chavez-Garzon, N., Pezoa, R., Suarez-Aldunate, P., \& Pilleux, M. (2025). Comparing automated valuation models for real estate assessment in the Santiago Metropolitan Region: A study on machine learning algorithms and hedonic pricing with spatial adjustments. \textit{PLOS ONE, 20}(3), e0318701. https://doi.org/10.1371/journal.pone.0318701

\bibitem[Urbania(2025)]{urbania2025}
Urbania. (2025). \textit{Urbania Index: índice de venta y alquiler, informe de mercado Lima, julio 2025}. https://urbania.pe/blog/urbania-index-lima/

\bibitem[VanderWeele y Ding(2017)]{vanderweele2017}
VanderWeele, T. J., \& Ding, P. (2017). Sensitivity analysis in observational research: Introducing the E-value. \textit{Annals of Internal Medicine, 167}(4), 268--274. https://doi.org/10.7326/M16-2607

\bibitem[Wager y Athey(2018)]{wager2018}
Wager, S., \& Athey, S. (2018). Estimation and inference of heterogeneous treatment effects using random forests. \textit{Journal of the American Statistical Association, 113}(523), 1228--1242. https://doi.org/10.1080/01621459.2017.1319839

\end{thebibliography}
